\section{Vấn đề và thất bại cụ thể của hệ thống}
\noindent \textbf{Hệ thống tham chiếu.} Nhóm lựa chọn hệ thống đặt lịch và hàng đợi khám ngoại trú trong OpenMRS 3~\cite{openmrs_demo} làm cơ sở so sánh. Các chức năng được xem xét gồm tiếp nhận yêu cầu khám, quản lý lịch hẹn, ghi nhận người bệnh đến khám và sắp xếp hàng đợi. Trên cơ sở đó, nhóm tập trung cải tiến việc đặt lịch và đề xuất bổ sung AI hỗ trợ tư vấn, phân loại mức độ ưu tiên trước khi khám.

\noindent \textbf{Vấn đề trọng tâm.} Việc phân bổ lịch hẹn cần tính đến người bệnh đến trực tiếp, đến muộn, đổi lịch hoặc bỏ hẹn. Nếu xử lý chưa phù hợp, một số thời điểm có thể xuất hiện hàng chờ dài trong khi các thời điểm khác còn trống lịch. Ngoài ra, việc thu thập triệu chứng và phân loại mức độ ưu tiên phụ thuộc vào nhân viên y tế, có thể gây chậm phản hồi khi nhiều yêu cầu tư vấn được gửi cùng lúc.

\noindent \textbf{Thất bại dịch vụ cụ thể.} Trong phạm vi đặt lịch và hỗ trợ tư vấn, phân loại, nhóm tập trung phân tích ba tình huống thất bại, bao gồm (1) \textit{Vắng mặt hoặc đến trễ:} Người bệnh đã xác nhận lịch hẹn nhưng không đến hoặc đến trễ, khiến khung giờ khám bị bỏ trống hoặc ảnh hưởng đến các lượt khám tiếp theo; (2) \textit{Chậm phản hồi tư vấn:} Nhiều yêu cầu được gửi cùng lúc khiến điều dưỡng hoặc bác sĩ chưa kịp xử lý, làm người bệnh phải chờ lâu để nhận được tư vấn; và (3) \textit{Phân loại mức độ ưu tiên chưa phù hợp:} Thông tin triệu chứng chưa đầy đủ hoặc cách đánh giá chưa thống nhất khiến người bệnh có thể được xếp mức độ ưu tiên chưa đúng, làm chậm việc tiếp nhận các trường hợp cần khám sớm.
